\setcounter{section}{18}

  \zwischenueberschrift{Übungsaufgaben}

\inputaufgabe
{}
{

Es seien $L,M,N$
\definitionsverweis {riemannsche Mannigfaltigkeiten}{}{.}
Zeige die folgenden Eigenschaften.
\aufzaehlungdrei{Die Identität
\maabbdisp {\operatorname{Id}_{  M  }} { M } { M
} {}
ist eine
\definitionsverweis {Isometrie}{}{.}
}{Wenn
\maabb {\varphi} { L } { M
} {}
eine Isometrie ist, so ist auch die
\definitionsverweis {Umkehrabbildung}{}{}
\maabb {\varphi^{-1}} { M } { L
} {}
eine Isometrie
}{Wenn
\maabb {\varphi} { L } { M
} {}
und
\maabb {\psi} { M } { N
} {}
Isometrien sind, so ist auch $\psi \circ \varphi$ eine Isometrie.
}

}
{} {}

\inputaufgabe
{}
{

Es sei $M$ eine
\definitionsverweis {riemannsche Mannigfaltigkeit}{}{.}
Zeige, dass die Menge der
\definitionsverweis {Isometrien}{}{}
auf $M$ eine
\definitionsverweis {Gruppe}{}{}
bilden.

}
{} {}

\inputaufgabe
{}
{

Es sei
\maabbeledisp {\varphi} {\R} {S^1 \subseteq \R^2
} { t } { \begin{pmatrix}  \cos  t  \\ \sin  t    \end{pmatrix}
} {.}
Zeige, dass $\varphi$ eine
\definitionsverweis {lokale Isometrie}{}{}
von
\definitionsverweis {riemannschen Mannigfaltigkeiten}{}{}
ist, wenn man $\R$ und $\R^2$ mit der
\definitionsverweis {euklidischen Struktur}{}{}
und $S^1$ mit der induzierten riemannschen Struktur versieht.

}
{} {}

\inputaufgabe
{}
{

Es sei
\maabb {\varphi} {\R^n } { \R^n
} {}
eine
\definitionsverweis {lineare Abbildung}{}{,}
der $\R^n$ sei mit seiner
\definitionsverweis {euklidischen Struktur}{}{}
und mit der induzierten
\definitionsverweis {riemannschen Struktur}{}{}
versehen. Zeige, dass $\varphi$ genau dann eine
\definitionsverweis {lineare Isometrie}{}{}
ist, wenn $\varphi$ eine
\definitionsverweis {Isometrie}{}{}
bezüglich der riemannschen Struktur ist.

}
{} {}

\inputaufgabe
{}
{

Es sei

\mavergleichskette
{\vergleichskette
{ I
}
{ \subseteq }{ \R
}
{  }{
}
{    }{
}
{   }{
}
}
{}{}{}
ein
\definitionsverweis {offenes Intervall}{}{}
und
\maabbdisp {\gamma} { I } { \R^n
} {}
eine injektive
\definitionsverweis {stetig differenzierbare Kurve}{}{.}
Wir fassen sowohl $I$ als auch $\R^n$ über ihre
\definitionsverweis {euklidische Struktur}{}{}
als
\definitionsverweis {riemannsche Mannigfaltigkeit}{}{}
auf. Es sei vorausgesetzt, dass

\mavergleichskette
{\vergleichskette
{ \gamma(I)
}
{ \subseteq }{ U
}
{ \subseteq }{ \R^n
}
{    }{
}
{   }{
}
}
{}{}{}
mit $U$
\definitionsverweis {offen}{}{}
eine
\definitionsverweis {abgeschlossene Untermannigfaltigkeit}{}{}
ist. Zeige, dass genau dann eine
\definitionsverweis {Isometrie}{}{}

\mavergleichskettedisp
{\vergleichskette
{ I
}
{ =} { \gamma(I)
}
{ \subseteq} { \R^n
}
{ } {
}
{ } {
}
}
{}{}{}
vorliegt, wenn  $\gamma$
\definitionsverweis {bogenparametrisiert}{}{}
ist.

}
{} {}

\inputaufgabe
{}
{

Es seien $L_1,L_2,M_1,M_2$
\definitionsverweis {riemannsche Mannigfaltigkeiten}{}{}
und seien
\maabb {\varphi_1} {L_1 } { M_1
} {}
und
\maabb {\varphi_2} {L_2 } { M_2
} {}
\definitionsverweis {Isometrien}{}{.}
Zeige, dass auch
\maabbdisp {\varphi_1 \times \varphi_2} {L_1 \times L_2 } { M_1 \times M_2
} {}
eine Isometrie ist.

}
{} {}

\inputaufgabe
{}
{

Wir betrachten den
\definitionsverweis {Diffeomorphismus}{}{}
\maabbeledisp {\varphi} {S^1 \times \R_+} { \Complex \setminus \{(0,0)\}  =  \R^2 \setminus \{(0,0)\}
} {(s,t) } { st
} {,}
wobei der Zylinder $S^1 \times \R_+$ mit der
\definitionsverweis {riemannschen Produktstruktur}{}{}
versehen sei. Beschreibe die riemannsche Struktur auf
\mathl{\R^2 \setminus \{(0,0)\}}{,} die sich ergibt, wenn $\varphi$ eine
\definitionsverweis {Isometrie}{}{}
werden soll.

}
{} {}

\inputaufgabegibtloesung
{}
{

Es seien
\mathkor {} {r} {und} {R} {}
positive reelle Zahlen mit

\mavergleichskette
{\vergleichskette
{ 0
}
{ < }{ r
}
{ < }{ R
}
{    }{
}
{   }{
}
}
{}{}{,}
wir betrachten den
\zusatzklammer {eingebetteten} {} {}
Torus

\mavergleichskettedisp
{\vergleichskette
{ T
}
{ =} { \left\{ (x,y,z) \in \R^3  \mid   \left(  \sqrt{x^2+y^2} -R  \right)^2 +z^2  = r^2         \right\}
}
{ \subseteq} { \R^3
}
{ } {
}
{ } {
}
}
{}{}{}
mit der
\definitionsverweis {induzierten riemannschen Struktur}{}{.}
Zeige, dass zu jedem

\mavergleichskette
{\vergleichskette
{ \alpha
}
{ \in }{ [0,2 \pi [
}
{  }{
}
{    }{
}
{   }{
}
}
{}{}{}
die Abbildung
\maabbeledisp {\Psi_\alpha} {\R^3} {\R^3
} { \left(  x   , \,   y  , \,  z                 \right) } { \left(  x \cos \alpha - y \sin  \alpha   , \,   x \sin \alpha + y \cos  \alpha  , \,  z                 \right)
} {,}
eine
\definitionsverweis {Isometrie}{}{}
von $T$ in sich induziert.

}
{} {}

\inputaufgabe
{}
{

Berechne die Länge des linearen Weges von
\mathl{\begin{pmatrix}  0  \\ 1   \end{pmatrix}}{} nach
\mathl{\begin{pmatrix}  1  \\ 1   \end{pmatrix}}{} in der
\definitionsverweis {Poincaréschen Halbebene}{}{.}

}
{} {}

\inputaufgabe
{}
{

Berechne den Flächeninhalt des Rechteckes in der
\definitionsverweis {Poincaréschen Halbebene}{}{,}
das durch die Eckpunkte
\mathl{(1,1),\, (5,1),\, (1,3),\, (5,3)}{} gegeben ist.

}
{} {}

\inputaufgabegibtloesung
{}
{

Zeige, dass durch die
\definitionsverweis {Abbildungen}{}{}
\maabbeledisp {\varphi} { {\mathbb H} } { U \left(  0,1  \right)
} { z} { { \frac{  z-  { \mathrm i}  }{ z + { \mathrm i}  } }
} {,}
und
\maabbeledisp {\psi} {  U \left(  0,1  \right)  } { {\mathbb H}
} {w} { { \frac{   (w+1)  { \mathrm i}  }{ 1-w  } }
} {,}
eine
\definitionsverweis {biholomorphe Abbildung}{}{}
zwischen der
\definitionsverweis {oberen Halbebene}{}{}
${\mathbb H}$ und der offenen Einheitskreisscheibe $U \left(  0,1  \right)$ gegeben ist.

}
{} {}

\inputaufgabe
{}
{

Bestimme die
\definitionsverweis {Ableitungen}{}{}
der
\definitionsverweis {Abbildungen}{}{}
\maabbeledisp {\varphi} { {\mathbb H} } { U \left(  0,1  \right)
} { z} { { \frac{  z-  { \mathrm i}  }{ z + { \mathrm i}  } }
} {,}
und
\maabbeledisp {\psi} {  U \left(  0,1  \right)  } { {\mathbb H}
} {w} { { \frac{   (w+1)  { \mathrm i}  }{ 1-w  } }
} {.}

}
{} {}

\inputaufgabegibtloesung
{}
{

Beschreibe die
\definitionsverweis {Abbildung}{}{}
\maabbeledisp {\psi} { \Complex \setminus \{1\} } { \Complex
} {w} { { \frac{   (w+1)  { \mathrm i}  }{ 1-w  } }
} {}
in reellen Koordinaten.

}
{} {}

\inputaufgabegibtloesung
{}
{

Bestimme die
\definitionsverweis {partiellen Ableitungen}{}{}
der
\definitionsverweis {Abbildung}{}{}
\maabbeledisp {\psi} { \R^2 \setminus \{  (1,0 )\} } { \R^2
} {  \begin{pmatrix} a \\b   \end{pmatrix} } {  { \frac{  1 }{ -2a+1+a^2+b^2 } }  \begin{pmatrix} -2b \\1-a^2-b^2   \end{pmatrix}
} {.}

}
{} {}

\inputaufgabe
{}
{

Bestimme die Durchstoßungspunkte zur Geraden, die durch
\mathl{\begin{pmatrix}  0  \\ 0\\  -1   \end{pmatrix}}{} und einem Punkt
\mathl{\begin{pmatrix}  a  \\ b \\  0   \end{pmatrix}}{} mit

\mavergleichskette
{\vergleichskette
{ a^2+b^2
}
{ < }{ 1
}
{  }{
}
{    }{
}
{   }{
}
}
{}{}{}
verläuft, mit dem zweischaligen Hyperboloid, siehe
[[Minkowskiraum/Dreidimensional/Zweischaliges Hyperboloid/Eingeschränkte riemannsche Struktur/Beispiel|Beispiel 18.11]].

}
{} {}

Es seien
\mathkor {} {L} {und} {M} {}
zwei
$C^k$-\definitionsverweis {Mannigfaltigkeiten}{}{.}
Eine
\definitionsverweis {stetige Abbildung}{}{}
\maabbdisp {\varphi} { L } { M
} {}
heißt ein
\definitionswortpraemath {C^k}{ lokaler Diffeomorphismus }{,}
wenn es zu jedem Punkt

\mavergleichskette
{\vergleichskette
{ P
}
{ \in }{ L
}
{  }{
}
{    }{
}
{   }{
}
}
{}{}{}
eine
\definitionsverweis {offene Umgebung}{}{}

\mavergleichskette
{\vergleichskette
{ P
}
{ \in }{ U
}
{ \subseteq }{ L
}
{    }{
}
{   }{
}
}
{}{}{}
gibt, dass $\varphi  \vert_U$ ein
\definitionsverweis {Diffeomorphismus}{}{}
zwischen $U$ und $\varphi(U)$ induziert.

  \zwischenueberschrift{Aufgaben zum Abgeben}

\inputaufgabe
{}
{

Wir betrachten die Ellipsoidoberfläche

\mavergleichskettedisp
{\vergleichskette
{ E
}
{ =} { \left\{ (x,y,z)  \mid   2x^2+3y^2+5z^2 = 1        \right\}
}
{ \subset} { \R^3
}
{ } {
}
{ } {
}
}
{}{}{}
mit der vom umgebenden Raum induzierten
\definitionsverweis {riemannschen Struktur}{}{.}
Bestimme die
\definitionsverweis {linearen Isometrien}{}{}
des $\R^3$, die eine
\definitionsverweis {Isometrie}{}{}
auf $E$ induzieren.

}
{} {}

\inputaufgabe
{4}
{

Es seien $L,M$
\definitionsverweis {differenzierbare Mannigfaltigkeiten}{}{,}
\maabbdisp {\varphi} { L } { M
} {}
sei ein
\definitionsverweis {lokaler Diffeomorphismus}{}{}
und $M$ sei eine
\definitionsverweis {riemannsche Mannigfaltigkeit}{}{.}
Zeige, dass es eine eindeutige riemannsche Struktur auf $L$ derart gibt, dass $\varphi$ zu einer
\definitionsverweis {lokalen Isometrie}{}{}
wird.

}
{} {}

\inputaufgabe
{3}
{

Man gebe ein Beispiel für eine
\definitionsverweis {riemannsche Mannigfaltigkeit}{}{}
$M$ und einen
\definitionsverweis {Diffeomorphismus}{}{}
\maabbdisp {\varphi} { M } { M
} {,}
der keine
\definitionsverweis {Isometrie}{}{}
ist.

}
{} {}

\inputaufgabe
{4}
{

Berechne die Länge des oberen Halbkreies mit Radius $\sqrt{2}$ von
\mathl{\begin{pmatrix}  1  \\ 1   \end{pmatrix}}{} nach
\mathl{\begin{pmatrix}  -1 \\ 1   \end{pmatrix}}{} in der
\definitionsverweis {Poincaréschen Halbebene}{}{.}

}
{} {}

\inputaufgabe
{4}
{

Berechne den Flächeninhalt des Kreies mit Mittelpunkt
\mathl{\begin{pmatrix}  2  \\ 2   \end{pmatrix}}{} und Radius $1$ in der
\definitionsverweis {Poincaréschen Halbebene}{}{.}

}
{} {}

\inputaufgabe
{4}
{

Beschreibe eine explizite
\definitionsverweis {Isometrie}{}{}
zwischen der
\definitionsverweis {Poincaréschen Halbebene}{}{}
und der oberen Schale des zweischaligen Hyperboloids.

}
{} {}

 [[Kategorie:Latexseite]]
